%% kapa: exam=1 round=28 year=2017 period=2 form=A subject=accounting subject_ko=회계학 seq=14 number=54
%% source: Q-Net 감정평가사 기출문제 (cst00309) artlSeq=5199603 file=extracted/1008/5199603/2017년 제28회 감정평가사 1차시험 문제지/2교시 A형.pdf
%% conversion: kordoc 4.24.51 -> tools/gen1.py (auto); review=proofread-sonnet; flags=none
\begin{question}{54}
상품매매기업인 (주)감평은 20x1년 1월 1일 특허권(내용연수 5년, 잔존가치 ￦0)과 상표권(비한정적 내용연수, 잔존가치 ￦0)을 각각 ￦100,000과 ￦200,000에 취득하였다. (주)감평은 무형자산에 대해 원가모형을 적용하며, 정액법에 의한 월할상각을 한다. 특허권과 상표권 회수가능액 자료가 다음과 같을 때, 20x2년도 포괄손익계산서에 인식할 당기비용은? (단, 20x2년 말 모든 무형자산의 회수가능액 감소는 손상징후에 해당된다.)
\par{\small\centering\begin{tabular}{|l|c|c|}
\hline
 & 특허권 & 상표권 \\ \hline
20x1년 말 회수가능액 & ￦90,000 & ￦200,000 \\ \hline
20x2년 말 회수가능액 & 35,000 & 120,000 \\ \hline
\end{tabular}\par}
\begin{choices}
  \item ￦45,000
  \item ￦105,000
  \item ￦120,000
  \item ￦125,000
  \item ￦145,000
\end{choices}
\end{question}
\answer{4}
